\subsection{Algebraic Elements and Minimal Polynomials}

\begin{definition} \label{def:algebra-evaluation-and-algebraic-element} \hlt{Evaluation and Algebraic Elements}\\
Let $\mathcal{A}$ be an $\mathbb{F}$-algebra and let $a\in\mathcal{A}$.
\begin{enumerate}[label=(\roman*)]
\setlength{\itemsep}{0pt}
\item For $p(x)=\sum\limits_{j=0}^{n}c_jx^j\in\mathbb{F}[x]$, set $p(a)\defeq\sum\limits_{j=0}^{n}c_ja^j$, where $a^0\defeq1_{\mathcal{A}}$. The \hlt{evaluation map at $a$} is $\operatorname{ev}_a:\mathbb{F}[x]\to\mathcal{A}$, $p\mapsto p(a)$.
\item The element $a$ is \hlt{algebraic over $\mathbb{F}$} iff $(\exists\,p\in\mathbb{F}[x]\setminus\{0\})\,p(a)=0$. Equivalently, $a$ is algebraic over $\mathbb{F}$ iff $\ker\operatorname{ev}_a\neq\{0\}$.
\end{enumerate}
\end{definition}

\begin{proposition} \label{prop:algebra-evaluation-homomorphism} \hlt{Evaluation Is an Algebra Homomorphism}\\
Let $\mathcal{A}$ be an $\mathbb{F}$-algebra and let $a\in\mathcal{A}$.
\begin{enumerate}[label=(\roman*)]
\setlength{\itemsep}{0pt}
\item The map $\operatorname{ev}_a:\mathbb{F}[x]\to\mathcal{A}$ is an $\mathbb{F}$-algebra homomorphism.
\item If $a$ is algebraic over $\mathbb{F}$, then $\ker\operatorname{ev}_a\triangleleft\mathbb{F}[x]$ is nonzero.
\end{enumerate}
\begin{proof}
\begin{enumerate}[label=(\roman*)]
\setlength{\itemsep}{0pt}
\item[$(i)$] Let $p(x)=\sum\limits_{i=0}^{m}c_ix^i$ and $q(x)=\sum\limits_{j=0}^{n}d_jx^j$. Bilinearity and scalar centrality imply $(pq)(a)=\sum\limits_{i=0}^{m}\sum\limits_{j=0}^{n}c_id_ja^{i+j}=p(a)q(a)$.\\
For $r\in\mathbb{F}$, the evaluation definition implies $(p+q)(a)=p(a)+q(a)$ and $(rp)(a)=rp(a)$.\\
Since $\operatorname{ev}_a(1)=1_{\mathcal{A}}$, $\operatorname{ev}_a$ is an $\mathbb{F}$-algebra homomorphism by Definition \ref{def:algebra-homomorphism}(i).
\item[$(ii)$] If $a$ is algebraic, then $\ker\operatorname{ev}_a\neq\{0\}$ by Definition \ref{def:algebra-evaluation-and-algebraic-element}(ii). $\ker\operatorname{ev}_a\triangleleft\mathbb{F}[x]$ by Proposition \ref{prop:algebra-kernel-first-isomorphism}(i).
\end{enumerate}
\hlt{Structural Commentary}: Multiplication verifies evaluation; its kernel is an ideal by kernel absorption.\\
\hlt{Key Insight}: Powers of one element commute, so evaluation preserves polynomial products in any algebra.
\end{proof}
\end{proposition}

\begin{definition} \label{def:algebraic-element-minimal-polynomial-algebra} \hlt{Minimal Polynomial of an Algebraic Element}\\
Let $\mathcal{A}$ be $\mathbb{F}$-algebra, $a\in\mathcal{A}$ algebraic. Ideal $\ker\operatorname{ev}_a$ nonzero by Proposition \ref{prop:algebra-evaluation-homomorphism}(ii). Has unique monic generator by Theorem \ref{thm:integers-and-polynomial-ring-are-pids}(iv).
\begin{enumerate}[label=(\roman*)]
\setlength{\itemsep}{0pt}
\item The \hlt{minimal polynomial of $a$} is this generator $m_a\in\mathbb{F}[x]$.\\
Equivalently, $m_a$ is the unique monic polynomial of least degree satisfying $m_a(a)=0$.
\item The \hlt{polynomial annihilator of $a$} is $\operatorname{ann}_{\mathbb{F}[x]}(a)\defeq\{f\in\mathbb{F}[x]:f(a)=0\}=\ker\operatorname{ev}_a=\langle m_a\rangle$.
\end{enumerate}
\end{definition}

\begin{theorem} \label{thm:finite-dimensional-algebras-algebraic-elements} \hlt{Finite-Dimensional Algebras Are Algebraic}\\
Let $\mathcal{A}$ be an $\mathbb{F}$-algebra with $\dim_{\mathbb{F}}\mathcal{A}<\infty$. Then $(\forall\,a\in\mathcal{A})\,a$ is algebraic over $\mathbb{F}$.
\begin{proof}
Let $n\defeq\dim_{\mathbb{F}}\mathcal{A}$ and let $a\in\mathcal{A}$. The $n+1$ vectors $1_{\mathcal{A}},a,\ldots,a^n$ are linearly dependent.\\
Thus $(\exists\,(c_0,\ldots,c_n)\in\mathbb{F}^{n+1}\setminus\{0\})\,\sum\limits_{j=0}^{n}c_ja^j=0$; $p(x)\defeq\sum\limits_{j=0}^{n}c_jx^j$ is nonzero, satisfies $p(a)=0$; $a$ algebraic by Definition \ref{def:algebra-evaluation-and-algebraic-element}(ii).\\
\hlt{Structural Commentary}: Linear dependence among initial powers constructs an annihilating polynomial for arbitrary $a$.\\
\hlt{Key Insight}: Any list of more than $\dim_{\mathbb{F}}\mathcal{A}$ powers must contain a nontrivial linear relation.
\end{proof}
\end{theorem}

\begin{example} \label{ex:algebraic-infinite-dimensional-algebra} \hlt{An Algebraic Infinite-Dimensional Algebra}\\
Let $\overline{\mathbb{Q}}_{\mathbb{C}}\defeq\{z\in\mathbb{C}:(\exists\,p\in\mathbb{Q}[x]\setminus\{0\})\,p(z)=0\}$. For $\alpha,\beta\in\overline{\mathbb{Q}}_{\mathbb{C}}$, choose nonzero $f,g\in\mathbb{Q}[x]$ with $f(\alpha)=g(\beta)=0$.\\
Set $m\defeq\deg f$ and $n\defeq\deg g$. Reduction by $f(\alpha)=g(\beta)=0$ proves $\mathbb{Q}[\alpha,\beta]=\operatorname{span}_{\mathbb{Q}}\{\alpha^i\beta^j:0\leq i<m,\ 0\leq j<n\}$.\\
For $r\in\mathbb{Q}$, $\alpha+\beta,\alpha\beta,r\alpha\in\overline{\mathbb{Q}}_{\mathbb{C}}$ by Theorem \ref{thm:finite-dimensional-algebras-algebraic-elements}. Since $1\in\overline{\mathbb{Q}}_{\mathbb{C}}$, this set is a $\mathbb{Q}$-algebra.\\
For $n\in\N$, the polynomial $x^n-2$ is Eisenstein at $2$ and hence irreducible over $\mathbb{Q}$.\\
Any root $\alpha_n\in\mathbb{C}$ has $[\mathbb{Q}(\alpha_n):\mathbb{Q}]=n$ and belongs to $\overline{\mathbb{Q}}_{\mathbb{C}}$. Thus $(\forall\,n\in\N)\,\dim_{\mathbb{Q}}\overline{\mathbb{Q}}_{\mathbb{C}}\geq n$.\\
Hence this algebra is infinite-dimensional although each of its elements is algebraic.
\end{example}

\begin{proposition} \label{prop:generated-algebra-by-one-algebraic-element} \hlt{The Algebra Generated by One Algebraic Element}\\
Let $\mathcal{A}$ be an $\mathbb{F}$-algebra, let $a\in\mathcal{A}$ be algebraic, and set $d\defeq\deg m_a$ and $\mathbb{F}[a]\defeq\langle\{a\}\rangle_{\operatorname{subalg}}$.
\begin{enumerate}[label=(\roman*)]
\setlength{\itemsep}{0pt}
\item $\mathbb{F}[a]=\operatorname{im}\operatorname{ev}_a=\{p(a):p\in\mathbb{F}[x]\}$.
\item $(\forall\,b\in\mathbb{F}[a])(\exists!\,r\in\mathbb{F}[x])\,[b=r(a)\land(r=0\lor\deg r<d)]$.
\item Set $\overline{\operatorname{ev}}_a:\mathbb{F}[x]/\langle m_a\rangle\to\mathbb{F}[a]$, $p+\langle m_a\rangle\mapsto p(a)$. Then $\overline{\operatorname{ev}}_a$ is an $\mathbb{F}$-algebra isomorphism.
\end{enumerate}
\begin{proof}
\begin{enumerate}[label=(\roman*)]
\setlength{\itemsep}{0pt}
\item[$(i)$] The identity $\mathbb{F}[a]=\operatorname{span}_{\mathbb{F}}\{1_{\mathcal{A}},a,a^2,\ldots\}$ holds by Proposition \ref{prop:generated-subalgebra-and-ideal-formulas}(i) with $X=\{a\}$.\\
This span equals $\{p(a):p\in\mathbb{F}[x]\}=\operatorname{im}\operatorname{ev}_a$ by Definition \ref{def:algebra-evaluation-and-algebraic-element}(i).
\item[$(ii)$] For $p\in\mathbb{F}[x]$, choose $q,r\in\mathbb{F}[x]$ by Theorem \ref{thm:polynomial-division-algorithm}(i) with $p=qm_a+r$ and $r=0$ or $\deg r<d$.\\
Since $m_a(a)=0$, $p(a)=r(a)$, so the reduced representation exists.\\
Suppose $r,s$ satisfy the degree condition and $r(a)=s(a)$. Then $r-s\in\ker\operatorname{ev}_a=\langle m_a\rangle$.\\
If $r\neq s$, then $\deg(r-s)<d$, contradicting $\deg(hm_a)\geq d$ for $0\neq h\in\mathbb{F}[x]$. Thus $r=s$.
\item[$(iii)$] $\ker\operatorname{ev}_a=\langle m_a\rangle$ by Definition \ref{def:algebraic-element-minimal-polynomial-algebra}(ii), and $\operatorname{im}\operatorname{ev}_a=\mathbb{F}[a]$ by (i). Asserted induced isomorphism holds by Proposition \ref{prop:algebra-kernel-first-isomorphism}(iii).
\end{enumerate}
\hlt{Structural Commentary}: Generated powers identify the image; division selects representatives; the kernel excludes duplicates.\\
\hlt{Key Insight}: Reduction modulo $m_a$ converts polynomial division into unique normal forms for elements of $\mathbb{F}[a]$.
\end{proof}
\end{proposition}

\begin{theorem} \label{thm:algebraic-element-annihilator-and-invertibility} \hlt{Invertibility from the Constant Term}\\
Let $\mathcal{A}$ be an $\mathbb{F}$-algebra, $a\in\mathcal{A}$ be algebraic, $m_a$ be its minimal polynomial. Then $a\in\mathcal{A}^{\times}\Longleftrightarrow m_a(0)\neq0$.
\begin{proof}
\begin{enumerate}[label=\roman*.]
\setlength{\itemsep}{0pt}
\item[$(\Rightarrow)$] Suppose $a\in\mathcal{A}^{\times}$ and $m_a(0)=0$. Then $m_a=xq$ for a nonzero $q\in\mathbb{F}[x]$ with $\deg q<\deg m_a$.\\
Equality $0=m_a(a)=aq(a)$ implies $q(a)=0$ by multiplication by $a^{-1}$, contradicting minimality of $\deg m_a$, so $m_a(0)\neq0$.
\item[$(\Leftarrow)$] Suppose $m_a(0)\neq0$, $m_a(x)=xq(x)+m_a(0)$. Then $0=m_a(a)=aq(a)+m_a(0)1_{\mathcal{A}}$, so $b\defeq-m_a(0)^{-1}q(a)$ satisfies $ab=1_{\mathcal{A}}$.\\
Since $q(a)\in\mathbb{F}[a]$, $aq(a)=q(a)a$; hence $ba=ab=1_{\mathcal{A}}$ and $a\in\mathcal{A}^{\times}$.
\end{enumerate}
\hlt{Structural Commentary}: Factoring handles a zero constant term; separating a nonzero constant constructs the inverse.\\
\hlt{Key Insight}: A nonzero constant term constructs $a^{-1}$; a zero term lowers the annihilator degree.
\end{proof}
\end{theorem}
